\documentclass[../main.tex]{subfiles}
\begin{document}

\textit{Chương này trình bày các phương pháp quản lý dự án mà hệ thống hỗ trợ, các thuật toán phía
sau chức năng phụ thuộc, và các công nghệ dùng để xây dựng hệ thống.}

\section{Quản lý dự án cho nhóm phần mềm nhỏ}
\subsection{Kanban}
Kanban trực quan hoá công việc dưới dạng thẻ di chuyển qua các cột biểu diễn trạng thái, và giới hạn
lượng việc đang làm để nhóm hoàn thành việc trước khi bắt đầu việc mới~\cite{anderson2010kanban}.
Ưu điểm là luồng công việc liên tục: không có vòng lặp cố định, và việc bị chặn hiện ra ngay. Hệ thống
dùng quy trình bốn trạng thái cố định:
\[
\text{TODO} \rightarrow \text{IN\_PROGRESS} \rightarrow \text{BLOCKED} \rightarrow \text{DONE} \rightarrow \text{TODO}.
\]
Vòng trạng thái này được cài đặt ở cả client và server và phải luôn giống nhau; task ở trạng thái
BLOCKED mang theo lý do, nên bảng công việc giải thích được \emph{vì sao} việc bị dừng.

\subsection{Scrum}
Scrum là một khung làm việc trong đó nhóm bàn giao công việc theo các vòng lặp có độ dài cố định gọi
là \emph{sprint}, thường từ một đến bốn tuần~\cite{schwaber2020scrum}. Các tạo tác gồm \emph{product
backlog} (mọi việc có thể làm), \emph{sprint backlog} (việc nhóm cam kết trong sprint này) và phần
tăng trưởng sản phẩm. Các sự kiện gồm sprint planning, daily scrum, sprint review và \emph{sprint
retrospective}---nơi nhóm nhìn lại cách làm việc và quyết định cần thay đổi gì. Hai thước đo chuẩn:
\begin{itemize}
  \item \textbf{Burndown}: khối lượng còn lại (story point) theo từng ngày của sprint, so với đường
        lý tưởng đi thẳng từ mức cam kết về 0.
  \item \textbf{Velocity}: số story point hoàn thành mỗi sprint, dùng để định lượng cam kết sprint
        sau cho thực tế.
\end{itemize}
Hệ thống hỗ trợ Scrum trên nền bảng Kanban: mỗi task hoặc nằm trong product backlog, hoặc thuộc một
sprint; story point chính là độ phức tạp của task (S=1, M=2, L=3, XL=5).

\section{Đồ thị phụ thuộc và đường găng}
\subsection{Sắp xếp topo}
Nếu task $A$ phụ thuộc task $B$ thì $B$ phải xong trước. Các task và phụ thuộc tạo thành đồ thị có
hướng $G=(V,E)$ với cạnh $B \rightarrow A$. Khi $G$ không có chu trình (\emph{đồ thị có hướng không
chu trình}, DAG), luôn tồn tại một \emph{thứ tự topo}---dãy mà mỗi task đứng sau mọi điều kiện tiên
quyết của nó---và có thể tìm bằng thuật toán Kahn~\cite{kahn1962topological}: lặp lại việc lấy một
task không còn điều kiện tiên quyết chưa xong, đưa ra, rồi xoá các cạnh đi ra của nó. Khi nhiều task
cùng sẵn sàng, hệ thống ưu tiên theo độ ưu tiên, rồi hạn chót, rồi thời điểm tạo, để thứ tự gợi ý đặt
việc khẩn nhất chưa bị chặn lên đầu. Chu trình (người dùng có thể vô tình tạo ra) không làm thuật
toán lỗi: các task còn lại trong chu trình được đưa xuống cuối.

\subsection{Đường găng}
Phương pháp đường găng (CPM)~\cite{kelley1959cpm} tìm chuỗi hoạt động phụ thuộc dài nhất; mọi chậm
trễ trên chuỗi này làm chậm cả dự án. Hệ thống tính chuỗi nặng nhất gồm các task \emph{chưa xong} bằng
tìm kiếm theo chiều sâu có ghi nhớ, với trọng số của task $t$ là
\[
w(t) = p(t)\cdot c(t), \qquad
p \in \{1,2,5,10\}\ \text{(LOW..CRITICAL)},\quad
c \in \{1,3,5,8\}\ \text{(S..XL)},
\]
và trọng số của chuỗi là tổng trọng số các task. Task trên đường găng được đánh dấu trên bảng và
trong màn hình DAG. Một lỗi tinh vi phát hiện trong dự án liên quan tới việc ghi nhớ khi gặp chu
trình: giá trị tính cho một lượt duyệt bị cắt ngang bởi chu trình phụ thuộc vào đường đi, nên lưu
cache khiến kết quả phụ thuộc thứ tự mảng đầu vào. Các kết quả như vậy giờ không bao giờ được cache.

\section{Phần mềm local-first}
Phần mềm local-first giữ bản chính của dữ liệu trên thiết bị người dùng và coi server là một đầu
đồng bộ~\cite{kleppmann2019localfirst}. Lợi ích là phản hồi tức thì và chạy được khi ngoại tuyến;
khó khăn chính là gộp các chỉnh sửa đồng thời. Hệ thống chọn chính sách thận trọng có chủ đích: trình
duyệt ghi mọi thay đổi vào IndexedDB trước và hiển thị ngay, sau đó mới đồng bộ với server; khi
ngoại tuyến, dự án \emph{cá nhân} (một thành viên) vẫn sửa được, còn dự án \emph{chung} chuyển sang
chỉ đọc, vì khi chưa có thuật toán gộp, hai người cùng sửa ngoại tuyến sẽ âm thầm ghi đè lên nhau.

\section{Mô hình ngôn ngữ lớn và đầu ra có cấu trúc}
Các mô hình ngôn ngữ lớn dựa trên kiến trúc Transformer~\cite{vaswani2017attention} có thể tóm tắt,
trích xuất và lập kế hoạch từ ngôn ngữ tự nhiên. Tuy nhiên đầu ra của chúng là văn bản tự do, và ứng
dụng không thể hành động dựa trên một đoạn văn. Vì vậy hệ thống yêu cầu mô hình trả JSON theo lược đồ cố
định (ví dụ \code{\{went\_well, to\_improve, action\_items\}}), loại bỏ khung Markdown, kiểm tra và
chuẩn hoá kết quả (giá trị enum lạ được đổi về mặc định; phụ thuộc trỏ tới task không tồn tại bị
loại bỏ), và chuyển sang câu trả lời tất định khi đầu ra không dùng được. Mô hình sử dụng là
\code{openai/gpt-oss-120b}~\cite{openai2025gptoss} chạy trên Groq qua API tương thích
OpenAI~\cite{groqdocs}.

\section{Công nghệ sử dụng}
\begin{table}[H]
\centering
\caption{Các công nghệ sử dụng}\label{tab:stack}
\renewcommand{\arraystretch}{1.25}
\begin{tabular}{|p{3.0cm}|p{4.2cm}|p{6.5cm}|}
\hline
\textbf{Tầng} & \textbf{Công nghệ} & \textbf{Mục đích} \\ \hline
Frontend & Vue 3, TypeScript, Pinia, Vite, Tailwind CSS & Ứng dụng một trang; component dùng Composition API và một store trung tâm~\cite{vuejs} \\ \hline
Lưu trữ cục bộ & IndexedDB (idb-keyval) & Bản sao local-first các task của từng dự án \\ \hline
Backend & Python, FastAPI, Pydantic v2 & REST API với lược đồ request/response có kiểu~\cite{fastapi} \\ \hline
Dữ liệu & SQLAlchemy 2, Alembic, MySQL 8.4 & ORM, migration có phiên bản, CSDL production (SQLite khi phát triển) \\ \hline
Bảo mật & JWT (HS256), bcrypt & Phiên không trạng thái~\cite{rfc7519jwt}, băm mật khẩu \\ \hline
AI & Groq API, Ollama (tuỳ chọn) & LLM trên cloud, LLM cục bộ tuỳ chọn \\ \hline
Kiểm thử & pytest, Vitest, Vue Test Utils, Playwright & Kiểm thử đơn vị, component, tự động hoá trình duyệt \\ \hline
Triển khai & Docker Compose, nginx, GitHub Actions, Cloudflare Tunnel & Container, reverse proxy, CI, truy cập công khai không cần mở cổng~\cite{cloudflaretunnel} \\ \hline
\end{tabular}
\end{table}

\section{Tài liệu như một hệ thống kiểm soát}
Tài liệu mô tả hệ thống sẽ sai ngay khi mã thay đổi mà tài liệu không đổi theo; dự án này kế thừa một
bản đặc tả mô tả những hành vi mà mã không hề có. Vì vậy nhóm giữ tám tài liệu cạnh mã nguồn (D1--D8,
liệt kê ở phần Mở đầu), với quy tắc: một thay đổi mã làm tài liệu sai thì chưa hoàn thành cho tới khi
tài liệu được cập nhật trong cùng commit. D7 ghi lại mọi quyết định quan trọng (\emph{DEC-nnn}) và
mọi lỗi phát hiện được (\emph{F-nn}) cùng lý do, còn D8 truy vết từng yêu cầu tới mã nguồn và test
kiểm chứng nó. Các số liệu ở Chương~4 lấy từ những bản ghi này.

\end{document}
